\begin{resumen}
La Evaluación de la Calidad de Acciones (\textit{Action Quality Assessment}, AQA) en video tiene aplicaciones en deporte, rehabilitación y educación, pero los modelos de referencia (I3D, SlowFast) presentan un coste computacional prohibitivo para dispositivos con recursos limitados. Esta tesis aborda dos preguntas relacionadas: ¿en qué medida un \textit{pipeline} liviano basado en arquitecturas modernas puede aproximarse al rendimiento de un Teacher 3D profundo en AQA?, y ¿qué componentes del \textit{pipeline} sostienen ese rendimiento? Se propone un \textit{pipeline} reproducible basado en TSM-MobileNetV2 y MobileNetV3, inicializados con pesos preentrenados en ImageNet y entrenados con AdamW, programación cosenoidal, precisión mixta y acumulación de gradientes. La evaluación cubre tres dominios (AQA-7 multi-disciplina deportiva, MTL-AQA clavados, JIGSAWS cirugía robótica) y dos Teachers de referencia (I3D-R50 e SlowFast-R50). En AQA-7, el \textit{pipeline} liviano alcanza correlaciones SRCC de $0.9021 \pm 0.005$ (TSM-MobileNetV2) y $0.8907 \pm 0.005$ (MobileNetV3) sobre tres semillas, frente a $0.9052$ del I3D y $0.9158$ del SlowFast-R50; la brecha se mantiene menor a $0.025$ SRCC y robusta entre semillas, utilizando apenas $9\%$ de los FLOPs y reduciendo la latencia a 40~ms por clip. Las ablaciones identifican el preentrenamiento ImageNet ($\sim 3$ puntos SRCC) y el módulo Temporal Shift ($\sim 1$ punto adicional) como los componentes responsables del rendimiento. La transferencia cruzada delimita el alcance: MTL-AQA $\to$ AQA-7 transfiere parcialmente (SRCC $\sim 0.55$); AQA-7 $\to$ JIGSAWS fracasa. Como análisis marginal, se evalúa el aporte de la destilación de conocimiento sobre el \textit{pipeline} liviano; sólo 1 de 6 configuraciones mejora con destilación, lo cual cuestiona la premisa heredada de la literatura de que la destilación es necesaria para cerrar la brecha entre Teacher y Student. El código fuente, las configuraciones experimentales y los \textit{checkpoints} del \textit{pipeline} se publican como recurso reproducible.
\end{resumen}
